\documentclass[12pt,letterpaper]{article}
\usepackage[T1]{fontenc}
\usepackage{mathptmx}
\renewcommand{\ttdefault}{lmtt}
\usepackage[letterpaper,margin=1in]{geometry}
\usepackage{amsmath,amssymb}
\usepackage{microtype}
\usepackage{booktabs,longtable,array}
\usepackage{enumitem}
\usepackage{url}
\usepackage[hidelinks]{hyperref}
\newcommand{\F}{\mathbb{F}}
\newcommand{\Z}{\mathbb{Z}}
\newcommand{\R}{\mathbb{R}}
\newcommand{\PP}{\mathbb{P}}
\DeclareMathOperator{\GL}{GL}
\newcolumntype{L}[1]{>{\raggedright\arraybackslash}p{#1}}
\setlength{\emergencystretch}{3em}

\title{Revision log for \emph{Toric codes over finite fields}\\[4pt]
\large Draft 1 of the corrected and expanded version}
\author{Prepared for David Joyner by Claude (Anthropic), November 22, 2026}
\date{}

\begin{document}
\maketitle

\section{Files and sources}

This log records the changes from the manuscript \texttt{toric\_code7.tex}
(dated September 1, 2003) to \texttt{toric-codes.tex} (Draft 1, November
22, 2026). The 2003 file is the version published as D.~Joyner, \emph{Appl.
Algebra Engrg. Comm. Comput.} \textbf{15} (2004), 63--79, arXiv:math/0208155.
The arXiv and journal data were checked online.

The reviser read three files in full. The first was the 2003 manuscript.
The second was the referee report (ChatGPT, Report Draft 1, November 22,
2026), file \path{toric_codes_referee_report.tex}. The third was the style
guide (Draft 2, November 11, 2026), file
\path{joyner-math-paper-style-guide.md}. The
original MAGMA and GAP programs \texttt{toric.mag} and \texttt{toric.g} were not
available. All computations were redone with new scripts.

\section{AI usage}

\begin{itemize}[leftmargin=*]
\item \textbf{OpenAI ChatGPT (GPT-6), referee.} It wrote the referee report on
the 2003 manuscript. The report identified the false conjectures, the false
injectivity lemma, the reversed distance inequality, the errors in the
simplex theorem, the decoding problems, the fan and area errors, and the
duplicated table row. It also asked for a minimum-distance certificate for
the $[49,11,28]_8$ code.
\item \textbf{Anthropic Claude (Claude Opus 5.5), reviser.} It checked every
claim of the report independently (Section~\ref{sec:referee}) and recomputed
every table. It wrote the revised paper and this log. The following material
is Claude's: the grid-lemma proofs of Hansen's (a)--(c), the exact form of
the simplex theorem, the duality and cyclic-code propositions, the repaired
decoding theorem with its toric corollary and worked example, the search
that found the decoding pairs, the analysis of the minimum words of the
principal code, and the graph-theoretic section. It also wrote all scripts.
\item \textbf{Verification levels} (style guide, Section~7). Level~1 means an
independent derivation or an exhaustive exact computation. Level~2 means a
check against a source actually consulted. Level~3 means reliance on memory.
Section~\ref{sec:levels} lists the levels of the main claims.
\item AI review is not an independent certification of the results.
\end{itemize}

\section{Response to the referee report}\label{sec:referee}

Each item gives the referee's point, the reviser's independent verdict, and
the action taken. ``Agree'' means that the point was rederived independently
(Level~1), not merely accepted.

\begin{longtable}{@{}L{0.27\textwidth}L{0.19\textwidth}L{0.46\textwidth}@{}}
\toprule
Referee item & Verdict & Action\\ \midrule
\endhead
3.1 First conjecture false & Agree (Level 1) & Counterexample recomputed: the
square code over $\F_9$ has $d=49<52$. The conjecture is withdrawn and given as
Example~6.6. Ruano (2007) had already published counterexamples, so the
literature review cites him.\\
3.2 Second conjecture false; $k\ge\dim H^0$ reversed & Agree (Level 1) & $d=56<58$
for $S_1$ over $\F_9$. The dimension inequality is corrected. Example~6.7.\\
3.3 Injectivity lemma false as quantified & Agree (Level 1) & Replaced by
Proposition~5.3, an exact criterion: $\dim C_U$ equals the number of residues
mod $q-1$. The ``$|a_i|<q$'' criterion is corrected to intervals of length
$\le q-2$ (Remark~5.4).\\
3.4 Decoding setup incompatible with its example & Agree (Level 1) & The
algorithm is restated for functions regular at the evaluation points
(Algorithm~7.1, Theorem~7.2). The hypotheses are now separation by $V_B$, as
in Pellikaan's error-correcting pairs. The toric criteria are proved in
Lemma~7.6 and Corollary~7.7. The old example is explained in Remark~7.5 and
replaced by Example~7.8, which decodes the $[49,11,28]_8$ code end to end
up to $5$ errors. A pair found by search raises the guaranteed radius to $6$.\\
4.1 Subcode distance inequality reversed & Agree (Level 1) &
Proposition~5.6 and the following paragraph.\\
4.2 Simplex theorem: items (1) and (5) false & Agree (Level 1) & Replaced by
Theorem~6.1 with exact codimensions and the exact distance
$(q-1)^{r-1}(q-1-a)$ of $C_{S_a}$, proved by the grid lemma (Lemma~4.2). This
is stronger than the referee's suggested puncturing bound. The $r=3$
counterexample is recomputed ($d=48<51$).\\
4.3 Upper bound needs a hypothesis & Agree, with one correction & Theorem~6.4.
The referee's repaired statement allows $1\le h\le q-1$. This is wrong at
$h=q-1$: the product of all $q-1$ factors vanishes on the whole torus and
gives the zero word. The theorem requires $h\le q-2$. The statement is also
extended to any primitive lattice direction.\\
4.4 Lattices, cycles and divisors misidentified & Agree & Section~3.2:
fans live in $N_\R$, characters in $M$, $P_1+\dots+P_n$ is a zero-cycle, and
$D_i$ is the closure of a codimension-one orbit. The old formula for $D_i$ is
removed.\\
4.5 Fan5 refinement still singular, fan6 area wrong, principal fan singular
& Agree (Level 1) & Fan5: the correct refinement inserts $(-1,0)$, and
$G=2aD_3+aD_4$ (Example~5.10). The 2003 text also had $d_3=a$ instead of $2a$.
Fan6: the area is $\tfrac12(d_1+d_2+d_3)^2$. The fan is $\PP^2$ in disguise,
so the whole table follows from Hansen (b) (Example~5.13). Principal fan:
determinants $24,6,6$. The paper explains that only the polytope matters
(Section~5.3 and Section~8.3).\\
4.6 Duplicated row $(5;2,3,3)$ & Agree (Level 1) & $k=14$; the row with
$k=15$ is removed. Hansen-(b) rows with $a\ge q-1$ are labeled as outside the
hypothesis.\\
4.7 Reed--Muller hypotheses inconsistent; ``subcodes'' imprecise & Agree &
Example~1.1 now uses $0\le\ell\le m(q-1)$ and $\ell=a(q-1)+b$ with $0\le
b<q-1$. The subcode statement is made precise in Theorem~6.1(3).\\
D1--D4 Local proof errors & Agree & Corrected in Propositions~7.3 and 7.4. The
contradiction is now with $[r,fh]=0$ for the chosen $f$. The vector $r$ has
length $n$, and indices are fixed. A complete example now exists
(Example~7.8). Complexity is stated as $O(n^3)$.\\
Section 7 of report: $[49,11,28]$ certificate & Done (Level 1) & $d\ge28$ by two independent
exhaustive computations: an information-set search, and a Gray-code
enumeration of all $8^{11}$ codewords giving the full weight distribution
(Table~1). The claim of being better than the tables is dated to 2002. Grassl's
tables (consulted November 22, 2026) now list $[49,11,29]_8$.\\
S1--S7 Exposition and style & Agree & See Section~\ref{sec:changes}.\\
\bottomrule
\end{longtable}

The referee did not recheck the numerical tables. All $55$ rows (including the duplicate and the principal example) were
recomputed here (Section~\ref{sec:computations}). Apart from the duplicated
row they all agree with the 2003 values, so the original computations were
sound.

\section{Changes from the 2003 manuscript}\label{sec:changes}

\begin{enumerate}[leftmargin=*]
\item \textbf{Front matter.} The title is kept. The subtitle is ``Draft 1.
Corrected and expanded version''. The date line reads ``Original draft:
2004. Corrected and expanded draft: 2026.'' The original publication is
identified in a \verb|\thanks| footnote, and the address is now
\texttt{wdjoyner@gmail.com}. The abstract opens with an AI-disclosure
sentence. The phrase ``efficient list decoding'' is replaced by an accurate
description, because the algorithm is not list decoding.
\item \textbf{Structure.} New sections: literature review (before and after
2004), background with examples, notation table, graph-theoretic connections,
summary with status table and open questions, glossary. The acknowledgments
now include the AI statement.
\item \textbf{Section 2 (now 4).} Notation is clarified, and the misprints
``isoceles'' and ``trapazoid'' are corrected. Hansen's theorem gets a status
label. New Lemma~4.2 and Proposition~4.3 prove (a)--(c). Part (d) is marked
as Remark~2026.1. The sentence ``The bounds on $q$ are best possible'' is
replaced by a proved statement about $k$. In Remark~4.4 the old value
$d_{RM}-2q+1=10$ (for $q=9$, $a=3$) is corrected to $37$.
\item \textbf{Section 3 (now 5).} The lattice and divisor conventions are
corrected (moved to Section~3). Lemma~3.3 (J.~Hansen) is replaced by
Proposition~5.3. New: Propositions~5.5 and 5.6 and the duality
Theorem~5.7. Example tables are recomputed. Proofs are given of the Cartier
criteria for fan1 and fan2. The ampleness criterion on $\F_m$ is corrected to
$d_2'+d_4'>0$ and $d_1'+d_3'>md_2'$. Fan identifications ($\PP^2/\mu_3$,
$\PP(1,1,m)$, $\F_m$, $\PP^2$) are added. The correspondence with Hansen's
type (d) is stated with $a$ and $b$ interchanged, as it must be.
\item \textbf{Section 4 (now 6).} The conjectures are withdrawn and replaced
by counterexamples (Examples~6.6 and 6.7) and proved bounds
(Theorems~6.1 and 6.4). Definition~4.4 (standard toric code) moved to
Definition~5.1. Example~4.7 has the corrected hypothesis $2c\le q-2$.
\item \textbf{Section 5 (now 7).} The decoding section is rebuilt as
described in Section~\ref{sec:referee}. The bullet list of alternative
decoders is kept, and its unclosed parenthesis is fixed.
\item \textbf{New Section 8.} Graph-theoretic connections. These are proved
results (Proposition~8.1), computational results (Example~8.2, the
resolution of the principal fan), and unverified observations (Remarks
2026.3--2026.6, Question~8.3).
\item \textbf{Misprints.} ``Be begin'', ``a a'', ``In other for'', ``Must of
this'', ``an new'', ``minumum'', ``it's'', ``torii'', ``$x+1$'' for $x_1$, and
``$P.G$'' are corrected.
\item \textbf{Typesetting.} Class \texttt{article} is kept, following the
style guide's default for an existing file. The paper uses 12pt Times, US
Letter and one-inch margins. \texttt{newtx} is not installed in the build
environment, so \texttt{mathptmx} is used, the style guide's fallback. All
cross-references and citations use \verb|\label|/\verb|\ref| and \verb|\cite|.
Citation keys follow the initials-plus-year rule, with \texttt{GAP00} and
\texttt{GUA02} for packages. The file compiles with \texttt{pdflatex} with no
errors, no undefined references, and no overfull or underfull boxes.
\end{enumerate}

\section{Original results versus results new in 2026}

\textbf{Original results, retained (corrected where needed):}
\begin{itemize}[leftmargin=*,itemsep=0pt]
\item the construction of toric codes from fans and divisors,
\item the examples and tables (Examples~5.8--5.13), including the
$[49,11,28]_8$ code,
\item the comparison with punctured Reed--Muller codes (Remark~4.4),
\item the inclusion statements (Proposition~5.6, Theorem~6.1(1)--(3)),
\item the segment upper bound (Theorem~6.4, corrected),
\item the error-locator algorithm and its two propositions
(Algorithm~7.1, Propositions~7.3 and 7.4, repaired).
\end{itemize}

\textbf{Withdrawn:}
\begin{itemize}[leftmargin=*,itemsep=0pt]
\item Conjectures 4.1 and 4.2 (false),
\item Lemma~3.3 (false as stated),
\item the distance inequality of Remark~3.2 (reversed),
\item items (1) and (5) of Theorem~4.5 (false),
\item the old decoding example (ill-defined).
\end{itemize}

\textbf{New in 2026:} Lemma~3.3 (orthogonality), Lemma~4.2 (grid lemma),
Proposition~4.3, Propositions~5.3, 5.5 and 5.6, Theorem~5.7, the exact
items (4) and (5) of Theorem~6.1, Theorem~7.2 in its present form,
Lemma~7.6, Corollary~7.7, Example~7.8, Proposition~8.1, Example~8.2, the
weight distribution (Table~1), the Gilbert--Varshamov comparison, and the
fan resolutions.

\section{Concordance of numbering}

\begin{longtable}{@{}ll|ll@{}}
\toprule
2003 & 2026 & 2003 & 2026\\ \midrule
\endhead
Eq.~(1) & (1.1) & Example 3.10 (fan6) & Example 5.13\\
Example 1.1 & Example 1.1 & Conjecture 4.1 & withdrawn; Example 6.6\\
Eq.~(2) & (4.1) & Conjecture 4.2 & withdrawn; Example 6.7\\
Theorem 2.1 & Theorem 4.1 & Remark 4.3 & Example 6.7 (text)\\
Remark 2.2 & Remark 4.4 & Definition 4.4 & Definition 5.1\\
Example 2.3 & Example 4.5 & Theorem 4.5 & Theorem 6.1\\
Remark 3.1 & Section 3.2 (text) & Remark 4.6 & Theorem 6.1 (hypothesis)\\
Remark 3.2 & Proposition 5.6 & Example 4.7 & Example 6.3\\
Eq.~(3) & (5.1) & Theorem 4.8 & Theorem 6.4\\
Eq.~(4) & (5.2) & Remark 4.9 & Remark 6.5\\
Eq.~(5) & (5.3) & Eq.~(7) & (7.1)\\
Lemma 3.3 & Proposition 5.3 & Proposition 5.1 & Proposition 7.3\\
Eq.~(6) & removed & Proposition 5.2 & Proposition 7.4\\
Example 3.4 (fan7) & Example 5.8 & Remark 5.3 & Theorem 7.2 (hyp.\ (C))\\
Example 3.5 (fan1) & Example 5.9 & Example 5.4 & Remark 7.5; Example 7.8\\
Example 3.6 (fan5) & Example 5.10 & & \\
Examples 3.7, 3.8 & Example 5.11 & & \\
Example 3.9 (fan2) & Example 5.12 & & \\
\bottomrule
\end{longtable}

\section{Computations}\label{sec:computations}

All scripts are in the folder \texttt{scripts/}. They are licensed MIT. The
LaTeX is CC BY 4.0. There is no git repository yet, so no commit hashes can be
given. They should be added when the files are committed.

\begin{longtable}{@{}L{0.34\textwidth}L{0.58\textwidth}@{}}
\toprule
Script & Purpose and result\\ \midrule
\endhead
\texttt{gf.py}, \texttt{toric.py} & Field arithmetic for $q\le27$ and
generator matrices of full-torus codes.\\
\texttt{mindist.c} & Exact minimum distance (information-set search with
disjoint information sets), weight distributions for $q=2^e$, and listing of
words of given weight. Compile: \texttt{gcc -O3 -o mindist mindist.c}.\\
\texttt{test\_mindist.py} & Agrees with brute force on 300 random codes over
$\F_q$, $q\le9$.\\
\texttt{verify\_tables.py} & Recomputes all 55 table rows (including the duplicate and the principal example). All agree with the
2003 values except the duplicated row (computed $k=14$). Output in
\texttt{verify\_tables.json}.\\
\texttt{wd\_principal.py} & Weight distribution of $[49,11,28]_8$ by
enumerating all $2^{33}$ $\F_2$-combinations. Minimum weight $28$,
$A_{28}=833$, total $8^{11}$.\\
\texttt{minwords\_principal.py}, \texttt{minwords\_explain.py} & The $119$
projective minimum words, their zero sets (three parallel lines of $AG(2,7)$),
and the $14$ exceptional trinomials.\\
\texttt{automorphisms\_\allowbreak principal.py} & Stabilizer of the exponent set in
$\GL_2(\Z/7)$: order $2$.\\
\texttt{verify\_referee.py} & All counterexamples, fan determinants and
weights, areas, the fan5 correction, and the Cartier and ampleness checks.\\
\texttt{verify\_hansen.py} & Hansen (a)--(d) against exact computation in 79
cases (35 of type (d)), all agreeing. The run stopped at its time limit
during $q=8$.\\
\texttt{gv\_check.py} & 21 of the 42 distinct parameter sets beat the finite
Gilbert--Varshamov bound.\\
\texttt{resolve\_fan.py} & Minimal resolutions and weighted cycle graphs.\\
\texttt{decode\_example.py} & Example~7.8: decodes all of 300 random trials
for each number of errors $t\le5$; succeeds in 50 of 300 trials at $t=6$ and
8 of 300 at $t=7$.\\
\texttt{decode\_simplex\_\allowbreak search.py} & The largest simplex $u+g(S_s)$ inside
$P_0$ has $s=6$, giving $t=5$.\\
\texttt{decode\_search.py}, \texttt{decode\_pair\_check.py} & Search over
convex shapes $A_0\subset[0,4]^2$ with maximal $B_0$, stopped before
completion. The best pair found has
$t=6$ ($|A_0|=7$, $d(C_{A_0})=28$, $|B_0|=15$, $d(C_{B_0}^\perp)=7$). Its
decoder succeeded in 200 of 200 trials with 6 errors and 33 of 200 with 7.\\
\bottomrule
\end{longtable}

\section{Verification levels of the main claims}\label{sec:levels}

\begin{longtable}{@{}L{0.5\textwidth}L{0.42\textwidth}@{}}
\toprule
Claim & Level\\ \midrule
\endhead
All results labeled ``proved here'' & 1 (derived in the paper)\\
All table entries, $[49,11,28]_8$, weight distribution, minimum words & 1
(exhaustive exact computation)\\
Hansen's Theorem 4.1(d) & 3 for the source statement; partial Level~1 by
computation (Remark 2026.1)\\
Publication data of Jo04, Ha00, Ha02, Ha13, Ha17, LS06, LS07, Ru07, Ru09, SS09, Li11,
Li17, BK13, CC12, Pe92, AF93, Al99, Ki03, BCP97 & 2 (bibliographic records and
abstracts consulted online)\\
Current tables entry for $[49,11]_8$ ($d\ge29$, $d\le34$) & 2 (codetables.de,
November 22, 2026)\\
Brouwer's 2007 table value $d\ge28$ & 2 (as reported on codetables.de)\\
Ruano (2007) gave counterexamples to the conjectures & 2 (abstract)\\
Section-level citations of Fulton~[Fu93] (\S2.1, \S2.5, \S2.6, \S3.4) & 3\\
Page citations kept from the 2003 text: Stichtenoth pp.~217--218,
MacWilliams--Sloane Ch.~13, Pless--Huffman p.~920, Ha02 \S2.3 & 3 (not
consulted)\\
Alon--F\"uredi theorem as stated in Section~8.2 & 3. The paper proves the
special case it needs (Lemma 4.2).\\
Baker--Norine page numbers; Skorobogatov--Vl\u{a}du\c{t} data;
Fitzgerald--Lax title; Brown--Kasprzyk~[BK13b] journal; years of
\texttt{toric.g} and of Joyner--Verrill & 3\\
Castryck--Cools gonality bound $\le$ lattice width & 2 (abstract and a
citing paper)\\
Seven-champions and $[36,19,12]_7$ results of Brown--Kasprzyk & 2
(abstracts)\\
\bottomrule
\end{longtable}

\section{Sources that could not be consulted}

The reviser had no access to the full texts of Fulton's book, Stichtenoth's
book, MacWilliams--Sloane, the Handbook of Coding Theory, J.~P.~Hansen's 2002
paper, S.~Hansen's theses, or Alon--F\"uredi. Their page and theorem numbers
in the paper are either the 2003 author's own citations, carried over, or
section-level citations from memory. Nothing in the paper's proofs depends on
them, except Hansen's Theorem 4.1(d) and the interpretive statements in
Section~8.3.

\section{Open items for David}

\begin{enumerate}[leftmargin=*]
\item Decide whether to keep the graph-theoretic Section~8 in the corrected
edition or move it to a companion note, in the spirit of the style guide's
rules for the 1981 and Satake papers.
\item Install \texttt{newtx} and switch from \texttt{mathptmx} if preferred.
\item Check the copyright position before posting this corrected version of
the published paper (style guide 8.6).
\item Commit the scripts to a repository and add the commit hash to
Section~9.
\end{enumerate}

\end{document}
